\documentclass[a4paper]{article}

\usepackage[english]{babel}
\usepackage[utf8]{inputenc}
\usepackage{amsmath}
\usepackage{graphicx}
\usepackage[colorinlistoftodos]{todonotes}
\usepackage{titling}

\pretitle{%
  \begin{center}
  \LARGE
  \includegraphics[width=1\textwidth]{portada4.png}\\[\bigskipamount]
}
\posttitle{\end{center}}

\title{}
\author{}
\date{}

\begin{document}

\maketitle

\tableofcontents
\clearpage

\section{Introduccion}
El problema a abordar va a ser la resolucion en matlab de la ecuacion:
\begin{equation}
f(x)=x+log(x)
\end{equation}
Es decir, queremos encontrar el punto $\overline{x}$ tal que $f($\overline{x}$)=0.

\section{Funciones usadas}
\subsection{PuntoFijo}
Se trata de la función puntofijo(f,xo,error,max1), donde:
\begin{itemize}
\item f es la función a hallar el punto fijo
\item xo es el punto inician donde se empieza la iteración
\item error como de preciso se requiere el resultadp
\item max1 es el número máximo de iteraciones que se quiere hacer
\end{itemize}
Ésta funcion lo que hace es ir iterando el punto la funcíon hasta que el error abosuto sea mas pequeño que el error pedido, a no ser de que se supere en número máximo de iteraciones, en cuyo caso no se habrá encontrado el punto buscado.

\section{Resolucion del Problema}
Primero definimos tres metodos iterativos:
\begin{equation}
\begin{array}{ccc}
\mathrm{I1:}x_{n+1}=log(x_{n}) & \mathrm{I2:}x_{n+1}=e^{-x_{n}} & \mathrm{I3:}x_{n+1}=\frac{x_{n}+e^{-x_{n}}}{2}
\end{array}
\end{equation}
Que surgen de despejar x de distintas formas. Ahora estas 3 iteraciones son iteraciones de punto fijo. Con lo cual primero tenemos que ver cuan de los tres es el mas optimo.

Primero vamos a ver donde estará aproximadamente la raíz de la ecuación. Como $f(0)<0$ y $f(1)>0$, entonces por Bolzano la raíz estará en el intervalo [0,1].

Sungundo, definimos como:
\begin{equation}
\begin{array}{ccc}
g_{1}:=log(x) & g_{2}:=e^{x} & g_{3}:=\frac{x+e^{-x}}{2}
\end{array}
\end{equation}
Ahora hay que ver cual de ellas tiene derivada menor que 1 para que el metodo del punto fijo funcione. Derivando cada uno tenemos que:
\begin{equation}
\begin{array}{ccc}
max\{|g_{1}'|\}_{[0,1]}=+\infty & max\{|g_{2}'|\}_{[0,1]}=0.36 & max\{|g_{3}'|\}_{[0,1]}=0.31
\end{array}
\end{equation}
Vemos que las funciones que cumplen la condición son $g_{3}$ y $g_{2}$ y como $g_{3}$ tiene menor valor,el más óptimo es $g_{3}$.

Por ultimo solo tenemos que poner en matlab que ejecute el metodo del punto fijo en el punto inicial 0 y nos da como resultado: 0.56714329.

\section{Conclusiones}
\section{Referencias}
[1] http://es.slideshare.net/emateucr/coloquios-ucr-2013gitsolis

\end{document}
